The multiplier controls how many candidates are refined after routing.
For the visited universe $U(q)$, JHQ defines \cite{jhq}
\begin{equation}
 c(\alpha)=\max\{k,\lceil\alpha k\rceil\},\qquad
 C_\alpha(q)=\operatorname{Top}^{\rm primary}_{\min\{c(\alpha),|U(q)|\}}U(q).
 \label{eq:6-budget}
\end{equation}
Probe selection determines $U$; the budget determines how much of it
survives primary screening. The analysis assumes $|C_\alpha|\geq k$;
insufficient candidates leave sentinel slots. At $k=10$, the tested
grid spans 20--2,000 slots without retraining the index.

\Cref{fig:5-query-path} highlights the budget-dependent work. The routed
candidates and number of primary scores stay fixed, but the scan's
admission and compaction work can change. A logical cost model at fixed
batch size and routing is
\begin{equation}
 T(\alpha)=T_{\rm base}+T_{\rm select}(c(\alpha))+
 T_{\rm refine}(c(\alpha))+T_{\rm final}(c(\alpha)).
 \label{eq:6-cost-model}
\end{equation}
Here $T_{\rm base}$ covers transfers, rotation, routing and primary-score
arithmetic; the other terms cover candidate selection, residual scoring
and final selection. Resource changes mean that smaller budgets need
not give shorter searches. The rule therefore chooses the smallest
acceptable budget, rather than the fastest measured one
(\cref{sec:7-budget-trace}). Each trial reruns a complete search,
including rotation, routing and primary scoring, even though it reuses
the trained index.

\paragraph{Position among adaptive ANN methods.}
The multiplier $\alpha$ expands the retained candidate set, like Faiss's
\texttt{k\_factor} and cuVS refinement \cite{faiss_refine,cuvs_refinement}.
Other controls use GT tuning \cite{faiss_autotune}, learned per-query
stopping \cite{li_termination} or predicted recall \cite{darth}. Here,
agreement with reference outputs selects one reusable budget for the
workload. Refinement still scores reconstructed vectors, whereas the
cuVS example refines with source vectors.
